\documentclass[10pt,a4paper]{article}

\usepackage[T1]{fontenc}
\usepackage[utf8]{inputenc}
\input{glyphtounicode}
\pdfgentounicode=1
\usepackage{charter}
\usepackage{anyfontsize}
\usepackage{microtype}
\usepackage[a4paper,left=1.7cm,right=1.7cm,top=1.3cm,bottom=1.1cm]{geometry}
\usepackage{xcolor}
\usepackage{enumitem}

\definecolor{accent}{RGB}{31,78,95}

\usepackage[colorlinks=true,urlcolor=accent,linkcolor=accent,%
  pdftitle={CV Tylden Hounsa},pdfauthor={Tylden Hounsa}]{hyperref}

\renewcommand\normalsize{\fontsize{10.5}{13}\selectfont}
\normalsize
\pagestyle{empty}
\setlength{\parindent}{0pt}
\setlength{\parskip}{0pt}

\setlist[itemize]{leftmargin=1.1em,labelsep=0.5em,itemsep=0.5pt,%
  topsep=1.5pt,parsep=0pt,partopsep=0pt,label={\color{accent}\textbullet}}

% Titre de section : capitales espacées, filet fin
\newcommand{\sect}[1]{%
  \vspace{9pt}\par
  {\color{accent}\small\bfseries\textls[120]{\MakeUppercase{#1}}}\par
  \vspace{-2.5pt}
  {\color{accent}\rule{\linewidth}{0.4pt}}\par
  \vspace{3.5pt}}

% Ligne d'entrée : gauche en gras, droite alignée
\newcommand{\entry}[2]{\textbf{#1}\hfill #2\par}

% Projet : titre, lien GitHub facultatif, technologies à droite
\newcommand{\proj}[3]{\textbf{#1}%
  \ifx\relax#2\relax\else\ \textcolor{accent}{\textbar}\ \href{#2}{GitHub}\fi
  \hfill{\small\itshape #3}\par}

\begin{document}
\raggedright

% En-tête
\begin{minipage}[t]{0.56\linewidth}
  {\fontsize{25}{27}\selectfont\bfseries Tylden HOUNSA}\par
  \vspace{12pt}
  {\color{accent}\bfseries Étudiant en L3 MIASHS}\par
  Recherche un stage dans la Data\par
\end{minipage}\hfill
\begin{minipage}[t]{0.42\linewidth}\raggedleft\small
  \href{mailto:tyldenhounsa@gmail.com}{tyldenhounsa@gmail.com}\\
  \href{tel:+33745639838}{07 45 63 98 38}\\
  Montpellier\\
  \href{https://www.linkedin.com/in/tylden-hounsa-b14559349}{Mon LinkedIn}\\
  \href{https://github.com/Ashborn539}{Mon GitHub}\\
  \href{https://tyldenhounsa.dev}{tyldenhounsa.dev}
\end{minipage}

\vspace{9pt}

Étudiant en troisième année de MIASHS à Montpellier. Projets menés de la collecte des données à la restitution, avec nettoyage, analyse statistique, visualisation et modèles de machine learning. Travail principalement en R et Python, SQL pour les bases de données.

\sect{Formation}
\entry{Université Paul-Valéry Montpellier 3 \textcolor{accent}{\textbar}\ Montpellier}{Depuis septembre 2024}
Licence 3 MIASHS : Mathématiques et Informatique Appliquées aux Sciences Humaines et Sociales\par

\sect{Projets}
\proj{Mortalité et pollution aux particules fines (2026)}{https://github.com/Ashborn539/projet-SDD-X-BDD}{R, SQL, ggplot2, R Markdown, WAMP}
\begin{itemize}
  \item Étude de la corrélation entre le taux de mortalité et le niveau de particules fines (PM2.5) selon les pays.
  \item Collecte, nettoyage, modélisation et visualisation des données.
  \item Conception de la base de données relationnelle et rapport complet rédigé en R Markdown.
\end{itemize}
\vspace{5pt}

\proj{Site web de vente de PC (en cours, 2026)}{https://github.com/Ashborn539/TyldenHounsa}{HTML, CSS, PHP, JavaScript, AJAX, jQuery, WAMP}
\begin{itemize}
  \item Interface utilisateur et fonctions d'achat en ligne. Logique applicative et interactions dynamiques gérées en PHP, AJAX et jQuery sur un serveur local WAMP.
\end{itemize}
\vspace{5pt}

\proj{Projet de machine learning sous Orange (2025)}{https://github.com/Ashborn539/ML_avec_orange}{Orange Data Mining}
\begin{itemize}
  \item Analyse d'un jeu de données réel : classification et régression (supervisé), clustering (non supervisé).
  \item Comparaison des modèles entre eux et interprétation des résultats.
\end{itemize}
\vspace{5pt}

\proj{Simulation proie-prédateur (2025)}{https://github.com/Ashborn539/ecosysteme-proie-predateur}{Python, Java, TCP/IP, sockets}
\begin{itemize}
  \item Modélisation d'un écosystème loups-moutons, simulation développée en Python.
  \item Interface graphique Java reliée en temps réel à la simulation par une architecture client-serveur TCP en local.
\end{itemize}
\vspace{5pt}

\proj{Bomberman, version terminale (2024)}{https://github.com/Ashborn539/terminal-bomberman}{Python, algorithmes}
\begin{itemize}
  \item Recréation du jeu en ligne de commande, avec gestion des états et de la logique de jeu sur des structures de données.
\end{itemize}

\sect{Compétences}
\renewcommand{\arraystretch}{1.2}
\begin{tabular}{@{}p{4.8cm}@{}p{\dimexpr\linewidth-4.8cm\relax}@{}}
  \textbf{Langages} & Python, R, SQL, Java, PHP, JavaScript, HTML/CSS \\
  \textbf{Analyse et visualisation} & pandas, NumPy, matplotlib, ggplot2, R Markdown, Excel, Tableau \\
  \textbf{Machine learning} & scikit-learn, Orange Data Mining, apprentissage supervisé et non supervisé \\
  \textbf{Outils} & Git/GitHub, VS Code, RStudio, WAMP/MySQL \\
\end{tabular}

\sect{Expérience professionnelle}
\entry{Pachamama \textcolor{accent}{\textbar}\ Chef de partie (temps partiel)}{Juvignac \textcolor{accent}{\textbar}\ Depuis août 2025}
\begin{itemize}
  \item Passage de commis de cuisine à chef de partie, avec la responsabilité complète du poste de production froid.
  \item Préparation des plats en service, rythme tenu sous pression en parallèle des études.
\end{itemize}
\vspace{5pt}

\entry{Café Jules \textcolor{accent}{\textbar}\ Commis de cuisine (saisonnier)}{La Grande Motte \textcolor{accent}{\textbar}\ De mai à septembre 2025}
\begin{itemize}
  \item Mise en place et préparation des plats en service, poste de production froid tenu en autonomie.
  \item Respect des exigences d'hygiène et de qualité.
\end{itemize}

\end{document}
